\section{Apéndice: problemas prácticos comunes}

Al entrenar RL profundo es inevitable encontrar problemas recurrentes; aquí se registran algunos casos mencionados en la clase o inferidos razonablemente:

\begin{enumerate}
    \item \textbf{Por qué la política presenta un sudden collapse}: esto suele ocurrir cuando el importance sampling weight es demasiado alto o el clipping supera 1.2. Se recomienda calcular `std(ratio)` y `max(ratio)` después de cada epoch de entrenamiento; si el promedio supera 0.12, revertir al checkpoint anterior.
    \item \textbf{Por qué existen trayectorias duplicadas en el replay buffer}: la causa suele ser un reset de alta frecuencia (el reset comienza desde estados cercanos). La solución consiste en insertar `prioritized eviction` en el buffer, eliminando primero los pares de transition con `similarity > 0.9`.
    \item \textbf{Cómo controlar el entropy decay}: en SAC, el coeficiente de entropy $\alpha$ no debe ser fijo; en su lugar, el value head debe estimar un entropy target (por ejemplo, 0.2) y ajustarlo automáticamente con `log_alpha`.
    \item \textbf{Cómo contrarrestar el reward hacking}: establecer `constraint auxiliary rewards`, por ejemplo una penalización de `-0.2` por llegar a estados no naturales, y usar el potential-based theorem en el reward shaping para garantizar que la optimal policy no cambie.
    \item \textbf{Cómo mantener el training trace}: usar structured logging para registrar `config, seed, buffer stats, loss curves`, y subir estos registros a un centralized experiment dashboard para facilitar el future debugging.
    \item \textbf{Cuándo se debe cambiar de algoritmo}: si la training variance supera la checkpoint window (por ejemplo, 10 episodes consecutivos con variance > 50\% of mean), se puede convertir el job en un PPO + SAC dual training, con PPO como anchor policy.
\end{enumerate}

\begin{warningbox}{No tomes el apéndice como un checklist}
Estos problemas son solo heuristics comunes y no sustituyen el ablation de un entorno específico. Incluso en un env similar, conviene verificar la variación longitudinal de `variance`, `kl` y `reward`, y ajustar los thresholds en consecuencia.
\end{warningbox}

\end{document}
